\documentclass{article}
\usepackage[T1]{fontenc}
\usepackage{lmodern}
\usepackage{graphicx}
\usepackage{amsmath,amssymb}
\usepackage[a4paper,margin=2cm]{geometry}
\usepackage[absolute,overlay]{textpos}
\setlength{\TPHorizModule}{1mm}
\setlength{\TPVertModule}{1mm}
\usepackage[indonesian]{babel}
\usepackage{hyperref}
\setlength{\emergencystretch}{2em}
% unit: Halaman 25

\begin{document}

% === Halaman 25 (llm) ===

\section*{Kurva Eliptik pada GF(p)}

\begin{itemize}
    \item Bentuk umum kurva eliptik pada GF(p) (atau $\text{F}_p$) :
\end{itemize}

\[ y^2 \equiv x^3 + ax + b \pmod{p} \]

\vspace{10mm}

yang dalam hal ini $p$ adalah bilangan prima dan elemen-elemen himpunan di dalam medan galois adalah $\{0, 1, 2, \dots, p-1\}$

\end{document}
